% ============================================================
%  SECCIÓN II: ANTECEDENTES — VERSIÓN EXPANDIDA
%  Estilo: IEEE Transactions
%  Paper: IA para Educación Inclusiva y Personalizada en CR
% ============================================================

\section{Antecedentes}

La problemática de la educación inclusiva en Costa Rica no es reciente ni
aislada; se inscribe en un patrón estructural documentado a lo largo de
sucesivos informes nacionales e internacionales. Esta sección organiza los
antecedentes en cuatro ejes: (A) el marco legal e institucional que obliga
al Estado a garantizar la inclusión educativa; (B) el diagnóstico actualizado
del sistema educativo costarricense; (C) los esfuerzos previos en tecnología
educativa y sus limitaciones; y (D) el estado del arte en IA aplicada a la
educación inclusiva y personalizada.

% ------------------------------------------------------------
\subsection{Marco Legal e Institucional para la Inclusión Educativa}
% ------------------------------------------------------------

El fundamento jurídico de la educación inclusiva en Costa Rica descansa en
tres instrumentos normativos complementarios. En primer lugar, la
\textit{Ley N\textsuperscript{o} 7600: Igualdad de Oportunidades para las
Personas con Discapacidad}~\cite{Ley7600} establece, desde 1996, la
obligación del Estado de garantizar el acceso, la permanencia y la promoción
de estudiantes con discapacidad en el sistema educativo regular, prohibiendo
cualquier forma de discriminación y exigiendo la provisión de apoyos y
adecuaciones curriculares. En segundo lugar, la \textit{Ley N\textsuperscript{o}
8661}~\cite{Ley8661}, que ratifica la Convención sobre los Derechos de las
Personas con Discapacidad de la ONU, refuerza el enfoque de derechos y sitúa
la inclusión como un imperativo de justicia social y no como una concesión
administrativa. En tercer lugar, la \textit{Ley N\textsuperscript{o} 8968 de
Protección de Datos Personales}~\cite{Ley8968} regula el tratamiento de
información sensible, lo que resulta directamente relevante para cualquier
sistema de IA que recopile datos sobre el estado emocional, el rendimiento
o las condiciones de aprendizaje del estudiantado.

A nivel de política educativa, el \textit{Consejo Superior de Educación}
aprobó en 2017 la política \textit{``La persona: centro del proceso educativo
y sujeto transformador de la sociedad''}~\cite{MEP2017politica}, que
reorienta el currículo costarricense hacia el desarrollo de competencias para
el siglo XXI, la atención a la diversidad y la incorporación ética de la
tecnología. Esta política constituye el marco pedagógico dentro del cual
cualquier intervención con IA debe inscribirse.

% ------------------------------------------------------------
\subsection{Diagnóstico del Sistema Educativo Costarricense}
% ------------------------------------------------------------

El \textit{Décimo Informe Estado de la Educación 2025}~\cite{PEN2025},
publicado por el Programa Estado de la Nación (PEN) del Consejo Nacional de
Rectores (CONARE), ofrece el diagnóstico más actualizado y comprehensivo del
sistema educativo costarricense. Sus hallazgos son alarmantes en tres
dimensiones que justifican directamente la presente investigación.

\textbf{Rezago en competencias básicas.} Los resultados de las pruebas PISA
2022 revelan que el estudiantado costarricense presenta niveles de comprensión
lectora y razonamiento matemático equivalentes a los de estudiantes de
primaria, con un retroceso de hasta diez años en comparación con los
estándares esperados para el nivel secundario~\cite{PEN2025}. Este rezago
es desproporcionadamente mayor en la población estudiantil con discapacidad
o necesidades educativas especiales, quienes enfrentan una doble barrera:
la brecha de aprendizaje preexistente y la insuficiencia de los apoyos
pedagógicos disponibles.

\textbf{Desarticulación institucional y ``pobreza de la evaluación''.}
El informe documenta una sucesión de decisiones fragmentadas que han
debilitado programas estratégicos sin ofrecer sustitutos efectivos. La
discontinuación en 2023 del \textit{Plan Integral de Nivelación Académica
(PINA)} sin evaluaciones técnicas previas, y el fracaso de la
\textit{``Ruta de la Educación''} —que nunca logró formalizarse en un
documento con presupuesto ni metas medibles— son ejemplos paradigmáticos
de esta desarticulación~\cite{PEN2025}. Paralelamente, la implementación
de pruebas nacionales estandarizadas de baja calidad técnica ha configurado
una ``pobreza de la evaluación'' que sobreestima las capacidades estudiantiles
y oculta el rezago real, impidiendo que el sistema genere información
accionable para la personalización del aprendizaje.

\textbf{Debilitamiento del acompañamiento docente en tecnología.}
El desmantelamiento del \textit{Programa Nacional de Informática Educativa
MEP-FOD (PRONIE-MEP-FOD)} ha dejado al sistema sin una política clara de
inclusión digital y acompañamiento pedagógico en tecnología~\cite{PEN2025}.
Este vacío es especialmente crítico en el contexto de las Aulas Integradas y
los Servicios de Apoyo Educativo, donde la disponibilidad de herramientas
tecnológicas adaptadas es más urgente. Adicionalmente, el informe reporta que
el 41\,\% del profesorado costarricense experimenta niveles significativos de
desgaste emocional, lo que reduce aún más la capacidad del docente de apoyo
para diseñar e implementar estrategias de personalización manual~\cite{PEN2025}.

En cuanto a la magnitud del problema en educación especial, las
\textit{Estadísticas de Servicios de Educación Especial 2014--2021} del
MEP~\cite{MEP2022stats} documentan que, en los Servicios de Apoyo Educativo
—que atienden a la gran mayoría de los estudiantes con discapacidad dentro del
sistema regular—, apenas el 55\,\% de los estudiantes matriculados se
encontraba alfabetizado en 2019, cifra que ascendió marginalmente al 57\,\%
en 2020. En los servicios de Atención Directa (Aulas Integradas), el
porcentaje de estudiantes alfabetizados no superaba el 35\,\% en el mismo
período, lo que evidencia la profundidad de la brecha de aprendizaje en la
población objetivo de esta investigación~\cite{MEP2022stats}.

% ------------------------------------------------------------
\subsection{Esfuerzos Previos en Tecnología Educativa y sus Limitaciones}
% ------------------------------------------------------------

Costa Rica cuenta con una tradición reconocida en la incorporación de
tecnología en la educación. Sin embargo, los modelos predominantes han
mostrado limitaciones estructurales que los hacen insuficientes para abordar
la personalización profunda requerida por el estudiantado con discapacidad.

El modelo histórico del \textit{PRONIE-MEP-FOD}, operativo desde la década
de 1990, se centró en la dotación de infraestructura (laboratorios de
cómputo) y en el uso de software educativo genérico. Aunque este programa
democratizó el acceso a la tecnología, su enfoque fue fundamentalmente pasivo
—consumo de contenido— y no adaptativo~\cite{PEN2025}. Los materiales
digitales no se ajustaban al ritmo, estilo de aprendizaje ni estado emocional
del estudiante, reproduciendo en el entorno digital las mismas limitaciones
de la instrucción presencial tradicional.

En paralelo, el MEP ha promovido la elaboración de \textit{Programas
Educativos Individualizados (PEI)} como mecanismo de personalización para
estudiantes con necesidades educativas especiales. Si bien el PEI representa
un avance conceptual significativo —al reconocer la necesidad de rutas de
aprendizaje diferenciadas—, su implementación práctica enfrenta la saturación
administrativa que documenta el \textit{Décimo Informe}~\cite{PEN2025}: los
docentes de apoyo dedican una proporción desproporcionada de su tiempo a la
elaboración y actualización de estos documentos, en detrimento del tiempo de
interacción directa con los estudiantes.

Más recientemente, el \textit{Plan MEP 2025--2026}~\cite{MEP2025plan}
contempla la integración de herramientas digitales y plataformas de
aprendizaje en línea como parte de la estrategia de recuperación de
aprendizajes post-pandemia. No obstante, este plan no especifica mecanismos
técnicos de personalización adaptativa para la población con discapacidad,
dejando un vacío operativo que la presente investigación busca abordar.

% ------------------------------------------------------------
\subsection{Política Nacional de IA y su Articulación con la Educación}
% ------------------------------------------------------------

En el plano de la política tecnológica, Costa Rica ha dado pasos concretos
hacia la institucionalización de la inteligencia artificial. La
\textit{Estrategia Nacional de Inteligencia Artificial 2024--2027
(ENIA)}~\cite{MICITT2024}, publicada por el Ministerio de Ciencia,
Innovación, Tecnología y Telecomunicaciones (MICITT), establece como
objetivo central promover el uso ético, seguro y responsable de la IA,
procurando beneficios para la ciudadanía y alineándose con las prioridades
nacionales de desarrollo sostenible. La estrategia reconoce explícitamente
la educación como uno de los sectores prioritarios para la adopción de IA,
e identifica la equidad e inclusión como principios transversales de
cualquier implementación~\cite{MICITT2024}.

El \textit{Plan de Acción ENIA 2025}~\cite{MICITT2025} operacionaliza estos
principios mediante líneas de acción específicas orientadas a la formación
de talento humano, el desarrollo de capacidades institucionales y la
promoción de un ecosistema de IA ético. Este plan proporciona el marco
normativo dentro del cual la integración de IA en las Aulas Integradas del
MEP debe inscribirse para garantizar su alineación con la política pública
nacional.

La publicación de la \textit{Guía para Docentes: Inteligencia Artificial en
clase}~\cite{MEP2026} por parte del MEP, en colaboración con UNESCO, las
Naciones Unidas y la Universidad Latinoamericana de Ciencia y Tecnología
(ULACIT), representa el primer esfuerzo sistemático por traducir la ENIA al
aula costarricense. La guía propone un marco de uso ético de la IA,
identifica herramientas recomendadas para distintos contextos pedagógicos y
subraya la importancia de la equidad, la accesibilidad y la preservación del
vínculo pedagógico humano~\cite{MEP2026}. Sin embargo, la guía no establece
protocolos técnicos específicos para la personalización adaptativa en
contextos de educación especial, lo que constituye precisamente la brecha
que esta investigación pretende llenar.

% ------------------------------------------------------------
\subsection{Estado del Arte en IA para Educación Inclusiva y Personalizada}
% ------------------------------------------------------------

A nivel internacional, la investigación sobre IA aplicada a la educación
inclusiva ha avanzado de manera acelerada en la última década, articulándose
en cuatro líneas temáticas que convergen en la propuesta de la presente
investigación: (i) sistemas adaptativos de aprendizaje para estudiantes con
discapacidad, (ii) detección del estado emocional y carga cognitiva en
entornos educativos mediados por IA, (iii) integración de la Taxonomía de
Bloom en sistemas de seguimiento cognitivo automatizado, y (iv) desafíos de
implementación en contextos de países en desarrollo y América Latina.

% - - - - - - - - - - - - - - - - - - - - - - - - - - - - - -
\subsubsection{Sistemas Adaptativos de IA para Estudiantes con Discapacidad}
% - - - - - - - - - - - - - - - - - - - - - - - - - - - - - -

El trabajo seminal de Ahmed et al.~\cite{Ahmed2025} constituye el referente
teórico central de esta propuesta. Los autores identifican un amplio espectro
de herramientas asistivas impulsadas por IA y aprendizaje automático (ML)
—desde lectores de pantalla y sistemas de reconocimiento de actividad humana
hasta robots humanoides para el apoyo de estudiantes con trastornos del
espectro autista— y proponen un \textit{marco de aprendizaje personalizado
en tiempo real} que utiliza algoritmos de aprendizaje reforzado para adaptar
dinámicamente los contenidos, el ritmo y la dificultad de las lecciones.
El sistema incorpora la \textit{Taxonomía de Bloom}~\cite{Bloom1956} como
eje de monitoreo cognitivo continuo y señala que los factores de estado
emocional, estilo de aprendizaje y progreso del estudiante son determinantes
para una personalización genuinamente inclusiva~\cite{Ahmed2025}.

Complementando este marco, Farhah et al.~\cite{Farhah2025} revisan la
integración de la IA generativa en sistemas de aprendizaje adaptativo y
detallan arquitecturas que combinan características de personalización y
accesibilidad para estudiantes con discapacidad. Su estudio enfatiza la
incorporación de señales de detección emocional y estrategias conductuales
adaptativas para ajustar la instrucción y el andamiaje cognitivo según las
necesidades individuales del estudiante, aportando criterios de ingeniería
y evaluación directamente aplicables al diseño de un prototipo para las
Aulas Integradas del MEP~\cite{Farhah2025}.

Desde una perspectiva de diseño de sistemas para educación especial,
Jadán-Guerrero et al.~\cite{Jadan2024} sintetizan evidencia sobre sistemas
tutores inteligentes y soluciones de computación afectiva desarrolladas
específicamente para este contexto. Los autores reportan patrones de diseño
y evidencia de casos que demuestran una mayor personalización del contenido
cuando se incorporan señales afectivas en tiempo real, al tiempo que
subrayan el rol indispensable del docente como mediador de la interacción
tecnológica~\cite{Jadan2024}. Por su parte, Bhimavarapu~\cite{Bhimavarapu2025}
implementa una red neuronal artificial bicapa (\textit{Bi-Stacked ANN})
alimentada con características demográficas, académicas y de autopercepción
de la discapacidad, logrando producir recomendaciones individualizadas de
acceso y apoyos de aprendizaje. La metodología de selección de características
mediante optimización por enjambre de partículas (PSO) ofrece un enfoque
replicable para perfilar las necesidades del estudiantado costarricense con
discapacidad y orientar políticas institucionales de inclusión~\cite{Bhimavarapu2025}.

% - - - - - - - - - - - - - - - - - - - - - - - - - - - - - -
\subsubsection{Detección del Estado Emocional y Carga Cognitiva en IA Educativa}
% - - - - - - - - - - - - - - - - - - - - - - - - - - - - - -

La dimensión afectiva del aprendizaje ha emergido como uno de los factores
más críticos para la efectividad de los sistemas tutores inteligentes (ITS),
particularmente en poblaciones con dificultades de aprendizaje. Dutt y
Ahuja~\cite{Dutt2024} desarrollaron intervenciones de tutoría inteligente
para estudiantes con dificultades de aprendizaje y midieron experimentalmente
las emociones inducidas y la carga cognitiva durante el uso del ITS. Sus
resultados vinculan directamente los estados emocionales positivos y la
reducción de la carga cognitiva con mejoras en la usabilidad y el compromiso
del estudiante, aportando evidencia empírica de que la adaptación emocional
no es un complemento opcional sino un componente funcional del sistema~\cite{Dutt2024}.

En la misma línea, el sistema AISTY~\cite{AISTY2024} combina estimación de
atención y afecto basada en visión computacional con componentes de IA
explicable (XAI) para adaptar el ritmo y la interacción en niños con trastorno
del espectro autista (TEA). El sistema monitorea la atención y los estados
emocionales inferidos a partir de señales visuales para ajustar las tareas
y los apoyos en tiempo real, mientras que los mecanismos de explicabilidad
permiten que padres y docentes comprendan las decisiones del sistema. Este
enfoque centrado en la transparencia algorítmica resulta especialmente
relevante para el contexto costarricense, donde la confianza institucional
y la participación de las familias son condicionantes clave para la adopción
de tecnología en el aula~\cite{AISTY2024}.

% - - - - - - - - - - - - - - - - - - - - - - - - - - - - - -
\subsubsection{Taxonomía de Bloom e IA: Seguimiento Cognitivo Automatizado}
% - - - - - - - - - - - - - - - - - - - - - - - - - - - - - -

La Taxonomía de Bloom~\cite{Bloom1956} ha experimentado una revitalización
significativa en el contexto de la educación mediada por IA, al proveer un
marco jerárquico que permite operacionalizar el progreso cognitivo del
estudiante en niveles medibles y accionables por algoritmos de adaptación.

Das et al.~\cite{Das2020} compararon un clasificador LDA etiquetado y un
clasificador de texto multiclase basado en BERT para asignar automáticamente
niveles cognitivos de Bloom a preguntas de evaluación, utilizando un conjunto
de más de 3.000 ítems y alcanzando precisiones del 83\,\% y 89\,\%
respectivamente. Esta capacidad de clasificación automática de la complejidad
cognitiva ofrece una base replicable para construir herramientas diagnósticas
alineadas con los indicadores de logro del currículo del MEP, permitiendo
que el sistema identifique en qué nivel de Bloom se encuentra cada estudiante
y ajuste el material de manera dinámica~\cite{Das2020}.

En el ámbito de la práctica docente, Barros et al.~\cite{Barros2025}
desarrollaron un estudio de investigación-acción en el que guiaron a docentes
en la planificación de lecciones secuenciadas según la Taxonomía de Bloom,
apoyadas por \textit{prompts} de IA generativa. Los resultados reportan un
fortalecimiento de la autonomía estudiantil y un mayor compromiso cognitivo
de orden superior en educación fundamental, y la metodología de desarrollo
profesional docente propuesta constituye un modelo operativo transferible al
contexto de formación de los docentes de apoyo del MEP~\cite{Barros2025}.
De manera complementaria, Kwan et al.~\cite{Kwan2025} examinan rediseños del
aula invertida (\textit{flipped classroom}) en los que la IA generativa
produce actividades previas a la clase alineadas con Bloom y soporta tareas
de orden superior durante la sesión presencial, incluyendo bucles de
retroalimentación formativa. El modelo de interacción estudiante-docente-IA
propuesto ofrece un esquema de integración directamente aplicable a los flujos
de trabajo pedagógico en las Aulas Integradas~\cite{Kwan2025}.

% - - - - - - - - - - - - - - - - - - - - - - - - - - - - - -
\subsubsection{IA para el Aprendizaje Personalizado en Países en Desarrollo
y América Latina}
% - - - - - - - - - - - - - - - - - - - - - - - - - - - - - -

La literatura sobre implementación de IA educativa en contextos de recursos
limitados y en América Latina es aún incipiente, pero creciente. Bravo
Condoy et al.~\cite{BravoCondoy2025} realizaron un estudio de métodos mixtos
en instituciones públicas ecuatorianas —contexto socioeconómico comparable
al costarricense— utilizando encuestas y entrevistas para mapear percepciones,
infraestructura y barreras de formación docente que afectan la adopción de
aprendizaje personalizado basado en IA. Sus resultados muestran interés
generalizado entre los actores del sistema, pero identifican la preparación
docente insuficiente y el acceso desigual a la tecnología como restricciones
primarias~\cite{BravoCondoy2025}. Esta evidencia regional es directamente
pertinente para anticipar y mitigar desafíos similares en el despliegue del
protocolo propuesto en escuelas costarricenses, especialmente en zonas rurales
con menor conectividad.

Desde una perspectiva global, el \textit{Informe de Seguimiento de la
Educación en el Mundo 2024} de la UNESCO~\cite{UNESCO2024} advierte que la
tecnología educativa —incluida la IA— corre el riesgo de ampliar las brechas
existentes si no se implementa con un diseño explícitamente inclusivo y con
mecanismos de gobernanza que protejan los derechos de los estudiantes más
vulnerables. El informe subraya que menos del 10\,\% de los países de
ingresos medios cuentan con marcos normativos específicos para el uso de IA
en educación especial, situando a Costa Rica —con su ENIA~\cite{MICITT2024}
y su Guía Docente IA~\cite{MEP2026}— en una posición de relativa vanguardia
regional, aunque con brechas de implementación técnica aún por cerrar.
Complementariamente, el \textit{Banco Mundial}~\cite{WorldBank2024} ha
identificado la personalización del aprendizaje mediada por tecnología como
una de las intervenciones con mayor potencial para acelerar la recuperación
de aprendizajes post-pandemia en la región, con énfasis en la necesidad de
soluciones de bajo costo y alta escalabilidad adaptadas a las realidades
institucionales de cada país.